% =====================================================================
\lezione{23/11/2026}{Verso l'NLP: Word2Vec e l'idea di DarkVec}
  {Prof.~Matteo Boffa}
  {perché l'NLP in un corso di cybersecurity; ripasso degli embedding e del self-supervised learning;
   NLP e bag of words (Lab~2); ipotesi distribuzionale; Word2Vec: CBOW e skip-gram, obiettivo, softmax,
   negative sampling, codice; modelli pre-addestrati; limiti; DarkVec: scanner, darknet, dati, ipotesi
   della co-occorrenza (slide~1--10)}
\label{lez:2026-11-23}
% =====================================================================

\begin{nota}[Fonte]
Trascrizione automatica della registrazione, revisionata, confrontata con le slide \emph{7-Word2Vec}
(tutte) e \emph{8-Darkvec} (fino alla slide~10, p.~49 del PDF).
\end{nota}

\section{Sintesi}
Prima lezione di teoria del Prof.~Boffa, che terrà le lezioni rimanenti, tutte sull'NLP. Il percorso parte da
Word2Vec (2013), il primo modo efficace di far ``capire'' a una rete le relazioni tra parole (il classico
\emph{king} $-$ \emph{man} $+$ \emph{woman} $\approx$ \emph{queen}), passa per un'applicazione di sicurezza
(DarkVec) e arriva ai transformer, base dei modelli linguistici attuali. Il messaggio della lezione: le
parole sono categorie senza label, e la supervisione si ricava dal \textbf{contesto} in cui compaiono (``you
shall know a word by the company it keeps''). Lo stesso vale per gli IP che contattano una darknet. Teoria
completa nei capitoli~\ref{cap:word2vec} e~\ref{cap:darkvec}.

\section{Avvisi pratici}
\begin{itemize}
  \item Calendario annunciato: martedì 24/11 lezione (si finisce DarkVec), mercoledì 25/11 i
        \textbf{transformer}. Dal 09/11 non ci sono state lezioni di teoria: 11/11 e 18/11 Lab~3
        (\rif{2026-11-11}{11/11}), 16/11 nessuna lezione (evento Cyber Monday).
  \item Su Word2Vec \textbf{non ci sarà un lab}: il codice mostrato serve solo a capire.
  \item In fondo alle slide di DarkVec ci sono una demo su GitHub e il link a un'altra lezione del docente
        che tratta DarkVec, agenti e, secondo il docente, \textbf{anticipa il Lab~4}.
  \item Il docente segue la chat per chi è collegato da casa: si può aprire il microfono e intervenire.
\end{itemize}

\begin{esame}[cosa chiede il docente su Word2Vec]
\begin{itemize}
  \item \textbf{Non} chiede di scrivere a memoria le formule. Una volta ha dato uno screenshot della formula
        (dello skip-gram o della sua probabilità) chiedendo da dove viene, cosa significa ogni termine e qual è
        l'intuizione. Si può essere imprecisi, purché l'intuizione sia giusta.
  \item Non confondere \textbf{non contestuale} e \textbf{contestuale}: Word2Vec è non contestuale (un vettore
        per parola), il transformer è contestuale. Secondo il docente è un errore frequente.
\end{itemize}
\end{esame}

\section{Contenuti}

\subsection{Perché l'NLP}
È un buon esempio di embedding (proiettare proprietà del mondo reale in uno spazio numerico che le conserva)
ed è il tema del momento: i modelli linguistici sono ovunque e preoccupano anche per le loro capacità
offensive in ambito cyber, tanto che esistono limiti e guardrail sull'uso dei modelli più potenti. L'obiettivo
del docente è arrivare a dare un'intuizione di come funziona un modello come GPT. Come apertura mostra un suo
lavoro del 2022: Word2Vec applicato alle sessioni bash raccolte da \textbf{honeypot} (servizi esca che
mostrano all'attaccante un servizio vulnerabile). Le sessioni, tutte diverse come stringhe, rivelano gruppi
di URL, percorsi e IP che \textbf{compaiono spesso insieme}.

\subsection{Ripasso: embedding e self-supervised learning}
Il one-hot è uno spazio ortogonale in cui tutto è equidistante (distanza $\sqrt{2}$) e non c'è nozione di
somiglianza; una rete, risolvendo un compito (es.\ protocolli sicuri o insicuri), impara uno spazio in cui la
distinzione emerge. Gli embedding si imparano solo con le reti neurali. Con la supervisione la rete può
imparare divisioni non volute (cane contro lupo che diventa erba contro neve); per questo oggi gli embedding
si ottengono soprattutto in modo \textbf{self-supervised}. Yann LeCun, interpellato sulle innovazioni più
importanti dell'ultimo decennio, ha messo al primo posto il self-supervised learning: senza, niente
transformer né modelli linguistici. Con le immagini l'idea di comprimere e ricostruire è intuitiva; con il
testo no: come capire che \emph{king} e \emph{queen} sono più vicini di \emph{king} e \emph{kangaroo}? E come
ottenere un embedding per una frase, o per una sequenza di IP che arriva a una darknet (scanner legittimo o
botnet?), senza label?

\subsection{NLP e bag of words}
NLP: comprendere, interpretare e \textbf{generare} linguaggio umano (il docente sottolinea la generazione).
Secondo i grafici dell'AI Index già nel 2023 i modelli raggiungevano le prestazioni umane in molti benchmark.
Prima delle reti c'era la \textbf{bag of words}: ogni frase è un sacchetto, ogni parola un oggetto; la
matrice ha una colonna per parola del vocabolario e una riga per frase. È ciò che si è fatto nel
\textbf{Lab~2} con le sequenze di API call, ``senza sapere di fare NLP''. Il vocabolario può essere noto
(dizionario inglese) o, come per i nomi delle API, da stimare (sul training). In \code{scikit-learn} è il
\code{CountVectorizer}. Ignora ordine e semantica (aprire un registro e aprire un file non hanno nulla in
comune), eppure funzionava bene.

\subsection{Word2Vec: l'idea}
Articolo del 2013, basato sull'\textbf{ipotesi distribuzionale} del linguista J.~R.~Firth (1957): ``you shall
know a word by the company it keeps''. Il significato di \emph{banking} viene dalle parole che lo circondano;
più esempi si hanno, più il significato è affidabile. Non conta tutto il contesto, solo le parole più vicine
(una finestra). L'intuizione: se l'embedding di \emph{banking} permette di predire le parole attorno (o
viceversa), è un buon embedding. Due varianti, scelte con un parametro nella libreria: \textbf{CBOW}
(contesto $\to$ parola centrale) e \textbf{skip-gram} (parola centrale $\to$ contesto), la più usata e quella
trattata.

\subsection{Skip-gram: obiettivo, probabilità, negative sampling}
Si massimizza la log-probabilità media delle parole nella finestra $[-c, c]$ attorno a ogni parola, per tutte
le frasi del corpus; $c$ è un iperparametro, di solito piccolo (più contesto, più calcolo). Graficamente: il
one-hot della parola centrale ($1 \times V$) moltiplicato per una matrice $V \times H$ dà l'embedding
($1 \times H$), uno spazio più piccolo che costringe la rete a tenere solo le dimensioni importanti. Una
seconda rete mappa le parole di contesto (positivi) e altre parole (negativi). Prima del training le
proiezioni sono casuali e non significano nulla.

La probabilità che una parola stia nel contesto della centrale è il prodotto scalare dei due embedding
normalizzato su tutto il vocabolario. Vettori paralleli portano lo stesso significato (come in teoria dei
segnali): massimizzare il prodotto scalare rende simili gli embedding delle parole che stanno insieme. È una
\textbf{softmax}: nell'articolo il problema è impostato come una classificazione. Il denominatore richiede
tutto il vocabolario (circa 171\,000 parole in inglese) a ogni predizione: non scala. Il \textbf{negative
sampling} tiene il contributo positivo e sostituisce il denominatore con $k$ negativi casuali; più $k$ è
grande, più ci si avvicina alla distribuzione vera, ma costa di più. Nel codice: due \code{nn.Embedding}
(centrale e contesto), score positivo e negativo; nella pratica si usa Gensim.

\subsection{Pre-training e limiti}
Per la prima volta non serve addestrare da zero per ogni compito: si pre-addestra una volta su un corpus
enorme senza label (es.\ Wikipedia) e si riusa, per esempio per suggerire la parola successiva sul telefono o
come base per un classificatore (\textbf{fine-tuning}: un IP mappato in un embedding e un classificatore sopra
che decide se è uno scanner benigno o Mirai). La stessa idea tornerà con i transformer. Limiti:
\begin{itemize}
  \item \textbf{non contestuale}: con corpora di biologia, favole e informatica, \emph{mouse} viene trascinato
        verso l'animale (due contributi contro uno) e non si adatta alla frase;
  \item \textbf{parole nuove}: ``petaloso'' non ha una riga nella matrice $V \times H$; aggiungere una parola
        significa allargare la matrice e riaddestrare (è il problema dei token del test assenti dal training
        visto nei lab, e di DarkVec con gli IP nuovi);
  \item \textbf{parole rare}: \emph{Sitzpinkler} ha una rappresentazione separata da \emph{sitzen} e
        \emph{pinkeln} e compare troppo poco per impararla. I modelli successivi lavorano a \textbf{token}
        (pezzi di parola: i limiti dei chatbot si contano in token), così la parola composta sfrutta le parti
        note;
  \item \textbf{dalle parole alle frasi}: i compiti utili sono su frasi (argomenti vietati, codice malevolo,
        traduzione, sentiment); con cinque parole si hanno cinque embedding e non è chiaro come combinarli.
\end{itemize}

\subsection{DarkVec: il problema}
Applicazione sviluppata nel gruppo di ricerca del docente (primo autore L.~Gioacchini). La domanda nasce
attorno al 2020 dopo la botnet \textbf{Mirai}, il primo attacco riuscito che ha usato dispositivi IoT
vulnerabili per campagne di denial of service. Gli scanner non hanno un nome, solo un IP che può cambiare:
bloccarlo dopo molti attacchi è tardi. Se so che un IP è di Mirai e vedo che le sue ``compagnie'' sono sempre
le stesse, posso dire qualcosa sugli altri IP prima di essere attaccato?

Per avere visibilità servono \textbf{vantage point}; il Politecnico ha darknet al Politecnico, in altre
università italiane, in Brasile, negli Stati Uniti e nel cloud. Una \textbf{darknet} occupa IP e porte senza
offrire servizi e senza rispondere: chi la contatta lo fa per scansione, attacco, errore di configurazione o
backscattering. È la parte ``a pensiero veloce'' dell'AI: volumi enormi e molto dinamici, che non lasciano il
tempo di usare un transformer. I dati sono solo IP e porte (nessuna richiesta segue, perché la darknet non
risponde), quindi \textbf{categorici} e \textbf{senza label} (tranne pochi IP segnalati); il monitoraggio
manuale si limita a filtri che arrivano tardi. Servono quindi embedding e self-supervision, e poi i downstream
task (IP più vicini, classificazione se ci sono dati).

\subsection{DarkVec: l'idea}
Una botnet è fatta di tanti host che arrivano \textbf{insieme} a fare cose simili: contattano lo stesso servizio
o compaiono insieme su vantage point o porte diverse. Come per le parole, ``si capisce cosa vuole un IP dalle
compagnie che frequenta''. Chi arriva per errore non compare con costanza insieme ad altri; gli scanner
benigni di solito pubblicano i propri IP; chi compare con costanza insieme ad altri e non è in elenco è
sospetto. La lezione si ferma qui (slide~10); il resto di DarkVec è previsto per il 24/11.

\section{Domande aperte e punti da verificare}
\begin{verifica}
\begin{itemize}
  \item Il docente data ``petaloso'' al 2015; la slide dice 2016 (fa fede la slide).
  \item Il racconto dell'onorificenza a LeCun nella trascrizione non è chiaro (anno e citazioni): il manuale
        riporta solo la risposta sul self-supervised learning mostrata in slide.
\end{itemize}
\end{verifica}
